\documentclass{report}
\usepackage{amsmath,amssymb}
\setlength{\parindent}{0mm}
\setlength{\parskip}{1em}
\begin{document}
\begin{center}
\rule{6in}{1pt} \
{\large Jakub Mare{\v c}ek \\
{\bf Iterative Methods for Integer Least Squares}}

School of Mathematics \\ The University of Edinburgh \\ James Clerk Maxwell Building \\ The King's Buildings \\ Edinburgh EH9 3JZ
\\
{\tt jakub.marecek@ed.ac.uk}\\
Danny Kershaw\end{center}

\newcommand{\R}{\ensuremath{\mathbbm{R}}}
\newcommand{\N}{\ensuremath{\mathbbm{N}}}
\newcommand{\prob}[2]{\underset{#1}{\rm Prob}\left[{#2}\right]}
\newcommand{\norm}[1]{\left\| {#1} \right\|}
\DeclareMathOperator{\st}{s. t.}
\DeclareMathOperator{\kron}{\,\otimes\,}
\DeclareMathOperator{\diag}{diag}

Least squares rank among the most fundamental problems in Scientific
Computing, Signal Processing, and Statistics.
Over the years, numerous variants of the problem have been studied. A
variant, which has attracted much attention in recent years, is:

{\sc Box-constrained Integer Least Squares (BILS)}: Given integers $m, n
> 0$, a subset of integers $Q$, matrix $H \in \R^{m \times n}$,
and $y \in \R^m$, return $x\in Q^n$ minimising $\norm{y - Hx}^2$.

The solving {\sc BILS} to optimality is hard, in a number of precise
senses. Notably, it is NP-hard to approximate within any constant factor
\cite{MR1462727}. We present an iterative solver for the problem, based
on novel convex programming relaxations.
Such relaxations strengthen the least squares relaxation, obtained by
omitting the integrality constraints. In each iteration, a small number
of violated valid inequalities are added, and the relaxation is
reoptimised.

There have been proposed numerous convex programming relaxations of {\sc
BILS}. (See Table~\ref{tab:relaxations} in the Appendix.)
They seem to be, however, either too expensive to compute (e.g.
high-dimensional semidefinite programming), or too weak (e.g. linear
programming). We propose a hierarchy of progressively stronger
relaxations in second order-cone programming (SOCP):
\begin{align}
\min c^T x \\
\st \left\| A_i x + b_i \right\|_2 & \leq c_i^T x + d_i, \quad i = 1,\dots,m \notag \\
Ax & = b, \notag
\end{align}
Notice that solving a SOCP is only moderately more expensive than solving
a linear program,
but allows one to express the least-squares objective.

Our relaxations are based on the reformulation attributed to Atamt{\" u}rk:
\begin{align}
& \min_{x \in Q^n} \norm{y - Hx} \label{eqn:ils} \\
= & \min_{x \in Q^n} x^T H^T H x - 2 y^T H x + y^T y \\
= & \min_{x \in Q^n, t \in \R^n} t^T t \st Hx - y \le t, y - Hx \le t \\
= & \min_{x \in Q^n, t_0 \in \R, t \in \R^n} t_0 \st |y - Hx| \le t, t^T
t \le t_0. \label{eqn:reformulated}
\end{align}
By removing the constraint $x \in Q^n$ in the original problem (\ref{eqn:ils}),
we would obtain the usual least squares relaxation.
By removing the same constraint in the reformulated problem (\ref{eqn:reformulated}),
we obtain the single-cone SOCP relaxation:
\begin{align}
\label{socp}
z_r = \min_{x \in R^n, t_0 \in \R, t \in \R^n } t_0 \st |Hx - y| \le t, t^T t \le t_0.
\end{align}
This relaxation (\ref{socp}) can be strengthened by the addition of
violated valid inequalities (``cuts''), including variants of
Chv{\'a}tal-Gomory cuts \cite{Cezik2005} and mixed-integer rounding (MIR)
cuts \cite{Atamturk2007}.

Let us define the Conic Mixed Integer Rounding (MIR) function.
For $0 \le f \le 1$, the conic rounding function $\phi_f : \R \to \R$ is:
$$\phi_f(a) = \begin{cases} (1 - 2f)n - (a - n) & \text{if $n \le a \le n + f$,} \\
(1 - 2f)n + (a - n) - 2f &\text{if $n + f \le a \le n + 1$.}
\end{cases}$$
The extension to vectors is element-wise.
For pure integer programs (\ref{socp}) and $\alpha \ne 1$, the conic
rounding cut hence is:
$$\sum_{j = 1}^{n} \phi_{f_{\alpha}}(a_j / \alpha) - \phi_{f_{\alpha}}(b
/ \alpha) \le t / |\alpha|.
$$
It is easy to see that for $\alpha \ne 1$ and $f_{\alpha} = b/\alpha -
\left\lfloor {b / \alpha} \right\rfloor$, conic rounding cuts are valid
for pure integer programs (\ref{socp}).
With each cut added, we obtain a new relaxation.

The choice of a cut has to take into account its violation and numerical
safety \cite{Cook2009}. After adding at most three cuts in a row, the
newly obtain relaxation is optimised, possibly using the previous
solution. The optimisation of the second-order cone program is, again,
performed using an iterative method, such as the primal-dual interior
point method.
The complexity of each iteration of the interior point method is
dominated the complexity of a matrix inversion of a $(m+c) \times (n+ m +
1)$ matrix, where $n,m$ are the dimensions of the original matrix $H$,
and $c$ is the number of cuts added. The appendix details promising
computational performance of this approach.

This approach lead to the development of low-power application specific
integrated circuits for the problem, which is of great importance in
mobile communications. There are systolic digital designs for much of the
interior point method already available.
This could be seen as an extension of Boyd's notion of real-time convex
programming to real-time integer convex programming.


\end{document}
